% Einheit 5 – Nr. 58–69
\setcounter{aufgabe}{57}
\einheitenkopf[e5][5 Potenzfunktionen]{Einheit 5 von 5 · Potenzfunktionen mit natürlichem Exponenten}
\zweigzeile{Hier lernst du Potenzfunktionen wie $y = x^3$ und $y = x^4$ kennen: Wertetabelle, Graph, Symmetrie und Funktionswerte · neu in diesem Jahr (am Gymnasium schon seit Klasse 9) · keine P10-Aufgabe · baut auf: Potenz mit dem Taschenrechner, Punkte eintragen, Normalparabel}

\verfahren{Wertetabelle und Graph}

\begin{aufgabe}{Ich kann eine Wertetabelle zu einer Potenzfunktion ausfüllen.}
Ergänze die Wertetabelle.
\begin{teile}
\teil $y = x^3$ \\ \wertetabelle{x}{y}{0,1,2,3}
\teil $y = x^4$ \\ \wertetabelle{x}{y}{0,1,2,3}
\teil $y = x^3$ \\ \wertetabelle{x}{y}{4,5,10}
\teil $y = x^4$ \\ \wertetabelle{x}{y}{4,5,10}
\teil $y = x^3$ \\ \wertetabelle{x}{y}{-3,-2,-1}
\teil $y = x^4$ \\ \wertetabelle{x}{y}{-3,-2,-1}
\end{teile}
\end{aufgabe}

\begin{aufgabe}{Ich kann eine Wertetabelle zu einer Potenzfunktion ausfüllen – weiter.}
Ergänze die Wertetabelle.
\begin{teile}
\teil $y = 0{,}5x^3$ \\ \wertetabelle{x}{y}{-2,-1,0,1,2}
\teil $y = -x^4$ \\ \wertetabelle{x}{y}{-2,-1,0,1,2}
\teil $y = x^3$ \\ \wertetabelle{x}{y}{{-1{,}5},{-0{,}5},{0{,}5},{1{,}5}}
\end{teile}
\end{aufgabe}

\begin{aufgabe}{Ich kann den Graphen einer Potenzfunktion zeichnen.}
Zeichne die Graphen mithilfe deiner Wertetabellen aus Nr. 58 und 59 in das Koordinatensystem und beschrifte sie.
\begin{teile}
\teil $y = x^3$
\teil $y = x^4$
\teil $y = 0{,}5x^3$
\end{teile}

\begin{ksys}[xmin=-2,xmax=2,ymin=-8,ymax=8,xstep=0.5,ystep=2]
\end{ksys}

\end{aufgabe}

\verfahren{Funktionswerte, Punktprobe und Umkehrung}

\begin{aufgabe}{Ich kann Funktionswerte einer Potenzfunktion berechnen.}
Es ist $f(x) = x^3$, $g(x) = x^4$ und $h(x) = 0{,}5x^3$. Berechne.
\begin{teilezwei}
\tz $f(2)$ \leerfeld & \tz $f(4)$ \leerfeld \\
\tz $g(3)$ \leerfeld & \tz $g(1)$ \leerfeld \\
\tz $f(-2)$ \leerfeld & \tz $g(-2)$ \leerfeld \\
\tz $h(-4)$ \leerfeld & \tz $f(1{,}5)$ \leerfeld \\
\tz $g(-0{,}5)$ \leerfeld & \\
\end{teilezwei}
\end{aufgabe}

\begin{aufgabe}{Ich kann prüfen, ob ein Punkt auf dem Graphen einer Potenzfunktion liegt.}
Prüfe mit der Punktprobe. Kreuze an.
\punktprobenliste{x^3/A(2|8), x^4/B(2|8), x^3/C({-3}|{-27}), x^4/D({-2}|{-16}), 0{,}5x^3/E(2|4), -x^4/F({-1}|1)}
\end{aufgabe}

\begin{aufgabe}{Ich kann zu einem Funktionswert das passende $x$ bestimmen.}
Bestimme alle $x$, die die Gleichung erfüllen.
\begin{teilezwei}
\tz $x^3 = 8$ \quad \feld{x} & \tz $x^3 = 1000$ \quad \feld{x} \\
\tz $x^3 = 1$ \quad \feld{x} & \tz $x^3 = 125$ \quad \feld{x} \\
\tz $x^4 = 16$ \quad \feld{x} & \tz $x^3 = -27$ \quad \feld{x} \\
\tz $x^4 = -1$ \quad \feld{x} & \\
\anweisung{Rechne mit dem Taschenrechner und runde auf zwei Nachkommastellen.}
\tz $x^3 = 50$ \quad \feld{x} & \\
\end{teilezwei}
\end{aufgabe}

\begin{aufgabe}{Ich kann am Funktionsterm erkennen, wie der Graph verläuft.}
Kreuze an, wie der Graph symmetrisch ist.
\begin{teile}
\teil $y = x^4$ \quad \kreuz{symmetrisch zur $y$-Achse}\kreuz{punktsymmetrisch zum Ursprung}
\teil $y = x^3$ \quad \kreuz{symmetrisch zur $y$-Achse}\kreuz{punktsymmetrisch zum Ursprung}
\teil $y = x^2$ \quad \kreuz{symmetrisch zur $y$-Achse}\kreuz{punktsymmetrisch zum Ursprung}
\teil $y = -2x^3$ \quad \kreuz{symmetrisch zur $y$-Achse}\kreuz{punktsymmetrisch zum Ursprung}
\end{teile}
\anweisung{Kreuze den Graphen an, der schmaler ist.}
\begin{teile}
\teil \kreuz{$y = 3x^4$}\quad\kreuz{$y = 0{,}5x^4$}
\end{teile}
\anweisung{Kreuze den Graphen an, der ganz unterhalb der $x$-Achse liegt (bis auf den Ursprung).}
\begin{teile}
\teil \kreuz{$y = x^4$}\quad\kreuz{$y = -x^4$}\quad\kreuz{$y = -x^3$}
\end{teile}
\end{aufgabe}

\begin{aufgabe}{Ich kann Gleichung und Graph einander zuordnen.}
Schreibe zu jeder Gleichung den Buchstaben des passenden Graphen.

\begin{ksys}[xmin=-2,xmax=2,ymin=-4,ymax=4,xstep=0.5,ablesen]
\funktion{\x^3}{}
\funktion{\x^4}{}
\funktion{-0.5*\x^4}{}
\funktion{0.25*\x^3}{}
\node[fill=white,inner sep=1pt] at (1.8,3.2) {$A$};
\node[fill=white,inner sep=1pt] at (-1.7,3.2) {$B$};
\node[fill=white,inner sep=1pt] at (1.85,-3.3) {$C$};
\node[fill=white,inner sep=1pt] at (1.75,2.2) {$D$};
\end{ksys}

\begin{teilezwei}
\tz $y = x^4$ \quad \leerfeld & \tz $y = x^3$ \quad \leerfeld \\
\tz $y = 0{,}25x^3$ \quad \leerfeld & \tz $y = -0{,}5x^4$ \quad \leerfeld \\
\end{teilezwei}
\end{aufgabe}

\begin{aufgabe}{Ich kann eine Potenzfunktion von einer Exponentialfunktion unterscheiden.}
Kreuze an. Rechne nichts.
\begin{teilezwei}
\tz $y = x^3$ \quad \kreuz{Potenz}\kreuz{Exponential} & \tz $y = 3^x$ \quad \kreuz{Potenz}\kreuz{Exponential} \\
\tz $y = 2x^4$ \quad \kreuz{Potenz}\kreuz{Exponential} & \tz $y = 4 \cdot 2^x$ \quad \kreuz{Potenz}\kreuz{Exponential} \\
\tz $y = x^2$ \quad \kreuz{Potenz}\kreuz{Exponential} & \tz $y = 0{,}5^x$ \quad \kreuz{Potenz}\kreuz{Exponential} \\
\end{teilezwei}
\end{aufgabe}

\begin{aufgabe}{Ich finde den Fehler bei Funktionswerten einer Potenzfunktion.}
Finde den Fehler, benenne ihn und korrigiere die Rechnung.
\begin{teile}
\teil Es ist $f(x) = -x^4$. Tim berechnet $f(2)$:
\rechnung{f(2) &= (-2)^4 \\ f(2) &= 16}
\teil Es ist $g(x) = x^3$. Sara berechnet $g(5)$:
\rechnung{g(5) &= 3 \cdot 5 \\ g(5) &= 15}
\end{teile}
\end{aufgabe}

\begin{aufgabe}{Ich kann begründen, wie der Graph einer Potenzfunktion verläuft.}
Begründe.
\begin{teile}
\teil Warum wird $y = x^4$ nie negativ?
\teil Warum verläuft der Graph von $y = x^3$ links unten durch den dritten Quadranten?
\end{teile}
\end{aufgabe}

\begin{aufgabe}{Ich kann mit einer Potenzfunktion eine Sachaufgabe zum Würfel lösen.}
Ein Würfel mit der Kantenlänge $k$ (in cm) hat das Volumen $V(k) = k^3$ (in cm³).
\begin{teile}
\teil Berechne das Volumen eines Würfels mit 5\,cm Kantenlänge.
\teil Ein würfelförmiges Gefäß fasst genau einen Liter (1000\,cm³). Wie lang ist seine Kante?
\teil Ein würfelförmiger Behälter soll 500\,cm³ fassen. Wie lang muss die Kante sein? Runde auf Millimeter.
\teil Wie ändert sich das Volumen, wenn man die Kantenlänge verdoppelt? Begründe.
\end{teile}
\end{aufgabe}
